\documentclass[12pt, a4paper]{article}

\usepackage[margin=2.5cm]{geometry}
\usepackage{times}
\usepackage{microtype}
\usepackage{parskip}
\usepackage{hyperref}
\hypersetup{colorlinks=false}

\begin{document}

% ── Title block ──────────────────────────────────────────────────────────────
\begin{center}
  {\Large\bfseries Enforcing Human Oversight in WireGuard VPN Deployments:\\[4pt]
  Design and Implementation of VISHWAAS}

  \vspace{10pt}

  {\normalsize
  Anurag Krishna Sharma,\quad HariOm Prakash,\quad Anitha Naidu,\quad J Jayachandran\\[2pt]
  \textit{NSED, Advanced Data Processing Research Institute}}

  \vspace{10pt}
  \noindent\rule{\linewidth}{0.4pt}
\end{center}

% ── Abstract ─────────────────────────────────────────────────────────────────
\begin{center}
  {\large\bfseries Abstract}
\end{center}

\noindent
Managing WireGuard VPN deployments across multiple machines presents a
persistent operational challenge. Static configuration files require manual
editing on every node whenever the network changes, making large or dynamic
deployments difficult to maintain. Automated peer-exchange tools such as
Tailscale and Nebula address the operational burden but do so by removing
human review from network access decisions entirely---nodes join and peer
automatically, leaving no clear record of who authorised a given connection
or when. In regulated or security-sensitive environments, this trade-off is
unacceptable.

This paper presents VISHWAAS (\textit{Sanskrit: trust}), a self-hosted
WireGuard VPN orchestration system designed around a single invariant: no
machine may join the network and no peer connection may become active without
explicit administrator approval. The system enforces this principle through a
two-tier architecture consisting of a central controller and a lightweight
per-node agent. The controller, built on FastAPI and React, exposes a web
dashboard through which an administrator reviews enrollment and connection
requests, approves or rejects them, and monitors the live state of the
network. Every decision is recorded in an append-only audit log with an event
type, actor, and millisecond-precision timestamp. The per-node agent translates
approved decisions into live WireGuard kernel configuration automatically,
eliminating the need for manual SSH access or configuration file editing.

The system provides three key-management strategies---agent-generated keys
(private key never leaves the node), controller-issued keys (centralised
generation with one-time delivery), and TPM~2.0 hardware-bound keys (private
key stored in Non-Volatile memory, never persisted to disk)---allowing
operators to select a security posture appropriate to their hardware and
threat model. Two network topology modes are supported: direct mesh for small
networks and hub-and-spoke gateway routing that reduces peer-count complexity
from $O(n^2)$ to $O(n)$ for larger deployments.

Security is enforced at two independent layers. The nginx reverse proxy
restricts all controller management endpoints to RFC-1918 private addresses
using a \texttt{geo} block, preventing internet exposure. The agent
additionally enforces per-request IP allowlisting at the application layer,
rejecting any request not originating from the registered controller address
before token validation even runs.

VISHWAAS was deployed on a three-machine internal lab network and validated
against five design goals: explicit approval for every state transition, a
complete audit trail, zero-touch provisioning after approval, flexible key
security escalating to TPM hardware binding, and confirmed absence of public
internet API exposure. The full enrollment-to-active-peer cycle completed in
under 30 seconds under normal conditions. The primary remaining limitations
are a single administrator account, no automated key rotation, and SQLite
persistence; all three constitute the immediate development roadmap and have
clear, well-understood implementation paths.

\noindent\rule{\linewidth}{0.4pt}

% ── Keywords ─────────────────────────────────────────────────────────────────
\noindent\textbf{Keywords:}\quad
WireGuard,\enspace VPN orchestration,\enspace human-in-the-loop,\enspace
network access control,\enspace TPM~2.0,\enspace audit logging,\enspace
FastAPI,\enspace self-hosted infrastructure

\end{document}
